\documentclass[11pt]{article}
\usepackage[utf8]{inputenc}
\usepackage[T1]{fontenc}
\usepackage{lmodern}
\usepackage[margin=1in]{geometry}
\usepackage{amsmath,amssymb}
\usepackage{textcomp}
\usepackage[dvipsnames]{xcolor}
\usepackage{enumitem}
\usepackage{parskip}
\usepackage{titlesec}
\usepackage[colorlinks=true,linkcolor=blue,citecolor=blue,urlcolor=blue]{hyperref}
\usepackage{natbib}
\usepackage{microtype}
\emergencystretch=1.5em

\titleformat{\section}{\large\bfseries}{\thesection}{1em}{}
\titleformat{\subsection}{\normalsize\bfseries}{\thesubsection}{1em}{}
\titleformat{\subsubsection}{\normalsize\bfseries\itshape}{}{0em}{}

\title{\textbf{The Certified-Human Anchor:\\ Temporal Holdout Validation for Corpus-Level AI-Content Estimators}}
\author{}
\date{}

\newcommand{\ph}{\textit{[placeholder --- fill in after experiments are run]}}

\begin{document}
\maketitle
\vspace{-2em}

\begin{quote}
\small\textit{Assembled from the six confirmed stage outputs of an \texttt{auto-research-discussion} run (seed topic B6.3), post research-taste audit (\textcolor{ForestGreen}{\textbf{PASS}}). Status: \textbf{specification / paper draft --- no experiments have been run.} All numbers below are pre-registered design targets, not results.}
\end{quote}

\begin{abstract}
Corpus-level AI-content estimators do not label individual documents; instead they report what fraction of a corpus is AI-generated by comparing it against a ``certified human'' reference distribution, almost always drawn from text that predates ChatGPT's late-2022 release on the assumption that such text cannot itself be AI-contaminated. We show this convention conflates two properties that come apart under time: a pre-2022 document's authorship can be genuinely human while the false-positive guarantee the anchor is assumed to provide is a property of the detector-distribution pair at a temporal distance from the certification date, not of the labels themselves --- so the guarantee can decay even when not one label is wrong. We construct a year-stratified corpus of hard-timestamped, certified pre-2022 human text, measure false-positive-rate drift across an existing battery of detectors and a mixture-based prevalence estimator, and disentangle any observed drift among three rival mechanisms --- genuine stylistic change, corpus-composition confound, and content leakage or estimator staleness --- using controls each immune to exactly one alternative. We further measure how much this drift propagates into a downstream corpus-level prevalence estimate under reference-year swaps, and define a validity envelope from a confidence bound rather than a raw significance threshold. We report where, if anywhere, the pre-2022 anchor's implied error guarantee holds, and where it does not.
\end{abstract}

\section{Introduction}

Estimating how much of a corpus is AI-generated has quietly shifted from a per-document classification problem to a population-level measurement problem. Instead of asking whether any single passage was written by a model, a growing line of work asks what fraction of an entire corpus --- a peer-review pool, a forum, a submission queue --- was touched by one, and answers it by fitting a mixture model against two reference distributions: one for AI-generated text, and one for genuinely human-authored text. That second reference is usually assembled from whatever text can be dated with confidence to before ChatGPT's public release in November 2022, on the reasoning that text written before large language models were widely accessible cannot itself be LLM-contaminated. This pre-2022 anchor is now a load-bearing, largely unexamined convention: nearly every corpus-level estimator that needs a human reference reaches for it.

Yet a pre-2022 document's authorship and the guarantee that document is assumed to license are not the same property, and they come apart under time. A document dated to 2015 can be genuinely, unambiguously human-written and still contribute to a false-positive rate that has drifted upward relative to whatever rate the anchor was calibrated to license, for reasons that have nothing to do with any label being wrong: human writing style itself is not static, and if it shifts across the calendar years a pre-2022 pool spans, a detector calibrated against recent pre-2022 text can flag older pre-2022 text at an elevated rate --- inflating an estimator's reported AI-prevalence, or misattributing ordinary stylistic change to AI generation, even though every document in the pool is exactly as certified-human as the convention assumes. This is not a hypothetical, and it is no longer entirely unmeasured: a recent audit of arXiv papers found the field's leading mixture estimator flags 13.1\% of pre-LLM-era review papers as AI-modified \citep{LLMGeneratedOrHumanWritten}, and a detector-adaptation study simulating a 2010 deployment found the same class of methods estimates 15.1\% AI content in 2020 human abstracts against only 2.3\% in 2010 \citep{MovingTargetTTA}. What these fragments establish is that the anchor's error is nonzero and year-dependent; what no estimator that depends on the anchor has done is measure that dependence as a function of temporal distance, attribute it among rival mechanisms, or bound where the anchor remains safe to use.

The nearest existing work approaches this gap from adjacent but distinct angles, and the closest of all is worth naming plainly before the wider landscape: the corpus-level mixture-estimation framework that popularized the pre-2022 anchor \citep{MonitorAIScale} fits a single mixing coefficient against a fixed human reference distribution and treats that reference as unproblematic ground truth, never testing whether its temporal composition confounds the estimated AI fraction --- it is the nearest prior work to this paper, and the one whose own instrument this paper directly audits, rather than merely gestures at. The one direct measurement of that reference's error to date \citep{LLMGeneratedOrHumanWritten} reports a pooled pre-LLM false-positive rate for this estimator and corrects for it as a static, type-specific offset, explicitly assuming the pattern is stable over time --- the assumption this paper tests rather than inherits. A longitudinal benchmark for concept drift in a different domain (malware detection) \citep{LAMDA} establishes that naive temporal evaluation inflates performance estimates and offers a template for year-stratified validation, but targets accuracy decay in a fixed feature space, not a calibrated false-positive bound for a text negative-anchor. And a cluster of work on contamination- and membership-detector fragility \citep{BenchContamFragility,MinKpp,ContamICL,MarkovCalib} shows that ``seen versus unseen'' separability collapses as two compared distributions converge or shift, but stops at that diagnosis without ever turning it into an explicit bound as a function of temporal distance --- the temporal gap between a pre-2022 anchor and a present-day corpus is exactly the kind of shift this literature describes in the abstract and has never measured concretely.

Turning this gap into a measurable question surfaces its own difficulties. An observed rise in false-positive rate over calendar years could stem from genuine stylistic change, from ordinary corpus-composition drift (older text skewing toward different genres, topics, or lengths than newer text), or from undetected AI content leaking into nominally pre-2022 buckets --- and a false-positive-rate curve alone cannot tell the three apart. Before that attribution question is even meaningful, the composition confound has to be ruled out: if year is confounded with genre in the available corpus, an apparent temporal trend may not be evidence about time at all. Compounding this, ``certified human'' is itself a metadata proxy rather than a verified property --- no oracle exists to certify that a document nominally dated before the cutoff involved no later revision, ghostwriting, or mislabeling, so any measured drift carries an unresolvable residual uncertainty unless that uncertainty is bounded rather than ignored. And because what actually matters to a downstream user is not a per-document detector statistic but the corpus-level prevalence number an estimator reports, a drift measured in raw false-positive rate says nothing on its own about whether it moves that number, since the mixing coefficient a prevalence estimator fits is a nonlinear function of the whole reference distribution, not a simple pass-through of per-document error.

Our approach turns each of these obstacles into a specific design choice rather than a hedge. Any observed drift is only attributed to genuine stylistic change once it survives two controls that are each immune to exactly one rival explanation --- a leakage-immune restriction to years unambiguously predating large-scale LLM availability, and a composition-matched resampling that equalizes topic and length distribution across year buckets --- with a synthetic-drift injection confirming, independently, that style-like perturbation alone can move the measured statistic in the same direction. The anchor's own residual impurity is bounded through a human-audited spot-check rather than assumed to be zero, folding that uncertainty into the reported bound instead of hiding it. Whether any of this matters downstream is tested directly, not inferred, by holding a target corpus fixed and swapping only the age of the human reference fed into an existing mixture-based prevalence estimator, reading off the resulting shift in the reported AI-fraction. And the resulting validity envelope is drawn from a confidence bound on the false-positive rate rather than a raw significance threshold, so that ``not enough data to tell'' and ``genuinely safe'' are never reported as the same thing.

We evaluate this design across multiple continuously-archived, hard-timestamped text sources spanning more than a decade of pre-2022 history, comparing the pre-2022 anchor's asserted guarantee against its measured behavior, alongside a year-shuffle placebo that checks the measurement pipeline itself is not manufacturing the very trend it is meant to detect.

In summary, we contribute:
\begin{itemize}[leftmargin=1.4em]
\item An identification and formalization of the asymmetry underlying the pre-2022 certified-human anchor: a document's authorship and the false-positive guarantee it is assumed to license are properties of different objects --- a label versus a detector-distribution pair at a temporal distance --- so the guarantee can decay even when every label is correct.
\item A disentanglement design that attributes any observed false-positive-rate drift among stylistic change, corpus-composition confound, and content leakage or estimator staleness, using controls each immune to exactly one alternative, rather than reporting an unattributed trend.
\item A direct measurement of whether and how much anchor-age drift propagates into the actual corpus-level number a downstream prevalence estimator reports, by swapping the human reference's age rather than inferring propagation analytically from document-level statistics.
\item A validity-envelope construction defined from a confidence bound rather than a raw significance threshold, reporting where the pre-2022 anchor's implied guarantee holds and where it does not, across multiple genres and a decade-plus time span.
\end{itemize}

\section{Problem Definition}

\noindent\textbf{Formal setup.} Let $y^*$ denote the certification cutoff (2022) and let $\Delta(y) = y^* - y$ denote the temporal distance of a calendar-year bucket $y$. A pre-2022 anchor corpus is partitioned into year-buckets $\{A_y\}$, each containing hard-timestamped, certified human-authored documents drawn from a fixed source. For an estimator $D_k$ operating at a fixed decision threshold, let $\mathrm{FPR}_k(y)$ denote its empirical false-positive rate on $A_y$ --- the fraction of certified-human documents in bucket $y$ that $D_k$ flags as AI-generated. The object of study is $\mathrm{FPR}_k(y)$ as a function of $\Delta(y)$: does it stay flat (the anchor's guarantee is uniform across the pre-2022 span), or does it grow with temporal distance (the guarantee erodes for older text)? A second object of study is the bias this induces in a downstream mixture-based prevalence estimator $\hat{\pi}(T; H)$, which estimates the AI-fraction of a fixed target corpus $T$ given a human reference $H$: does swapping $H$ across year-buckets shift $\hat{\pi}(T; H)$ --- that is, does drift measured at the document level actually move the number a downstream user consumes?

\noindent\textbf{Why this is hard.} Any relationship between $\mathrm{FPR}_k(y)$ and $\Delta(y)$ is compatible with at least three different mechanisms --- genuine secular change in human writing style, a corpus-composition confound (older buckets skewing toward different genres, topics, or lengths than newer ones), and undetected AI-content leakage into nominally pre-cutoff buckets --- and the curve alone cannot distinguish them. Before any of these can even be evaluated, the composition confound must be ruled out: if $y$ and corpus composition are entangled, a change in $\mathrm{FPR}_k(y)$ need not be a temporal signal at all. Compounding this, the label ``certified human'' that defines each $A_y$ is a metadata proxy --- a publication-year timestamp --- rather than a verified property of authorship, so any residual mislabeling contributes an unresolvable source of noise unless it is explicitly bounded. Finally, $\mathrm{FPR}_k(y)$ is a document-level statistic, while $\hat{\pi}(T; H)$ is a nonlinear function of the entire reference distribution $H$; a nonzero drift in $\mathrm{FPR}_k(y)$ does not by itself imply any particular shift in $\hat{\pi}(T; H)$, so the practical stakes of the drift can only be established by intervening on $H$ directly, not by extrapolating from the document-level curve.

\section{Method Overview}

At a conceptual level, the study is a single audit pipeline with five cooperating parts that turn a suspected but unmeasured anchor weakness into a bounded, attributed, and propagated quantity. Year-stratified anchor text and an existing battery of detectors and a mixture-based prevalence estimator enter the pipeline; a pair of disentanglement controls turns a raw false-positive-rate trend into an attributed claim about which of three rival mechanisms produced it; a purity-bounding step turns the anchor's own unverifiable authorship claim into an explicit uncertainty rather than a silent assumption; a reference-swap procedure turns a document-level drift signal into a measured shift in the corpus-level number a downstream estimator actually reports; and a confidence-bound envelope construction turns the fitted drift curve into a stated boundary condition rather than an all-or-nothing verdict.

The \textbf{disentanglement controls} take the raw false-positive-rate series for each year bucket and each detector and test it against two independent, mechanism-specific alternatives --- a leakage-immune sub-range and a composition-matched resample --- plus a synthetic-drift injection that checks whether style-like perturbation alone can reproduce the observed direction. The \textbf{purity-bounding procedure} takes a small human-audited sample of the anchor sources and turns it into an explicit residual-impurity estimate that is reported alongside, not folded silently into, every downstream bound. The \textbf{composition-matching procedure} takes the raw year-bucket samples and produces a matched resample in which topic and length distributions no longer vary with year, isolating whether the drift signal is genuinely temporal. The \textbf{reference-swap propagation test} takes a fixed target corpus and the mixture-based prevalence estimator and measures how much the estimated AI-fraction moves when only the age of the human reference changes. The \textbf{confidence-bound envelope construction} takes the fitted drift curve and produces the largest contiguous year-range whose upper confidence bound on false-positive rate stays within a pre-declared tolerance of the nominal calibrated rate, distinguishing a genuinely certified boundary from one that is merely under-sampled.

\section{Methodology}

\subsection{Disentanglement Controls}
\textbf{Challenge.} An observed false-positive-rate slope over calendar years is compatible with three different mechanisms --- genuine stylistic drift, undetected AI-content leakage into nominally pre-2022 buckets, or staleness of the estimator itself --- and no single statistic separates them.

\textbf{Insight.} No one control rules out all three at once, but each control can be built to be blind to exactly one alternative, so requiring survival across independently-designed controls narrows the explanation by elimination rather than by assumption.

\textbf{Method.} The drift regression is refit restricted to years unambiguously predating large-scale LLM availability (a leakage-immune sub-range), where any positive slope cannot be explained by leaked AI content by construction. Separately, graded lexical and syntactic perturbations are injected into a single fixed year's human sample at increasing intensity, and the flag rate is checked for a monotone rise with intensity --- establishing, independently of any real year-over-year data, that style-like change alone can move the measured statistic in the observed direction. A slope's sign is required to replicate across a majority of independently-built detector families before it is attributed to genuine drift rather than any one estimator's idiosyncrasy.

\subsection{Anchor Purity Bounding}
\textbf{Challenge.} ``Certified human'' is a metadata proxy --- a publication-year timestamp --- not a verified property of authorship, and no scalable oracle exists to certify that a nominally pre-2022 document involved no later revision, ghostwriting, or mislabeling.

\textbf{Insight.} Purity cannot be certified exactly at corpus scale, but the uncertainty it contributes can be bounded and folded into the reported result instead of silently assumed to be zero --- the same bound-rather-than-assume move this paper applies to the anchor's temporal validity, applied recursively to the anchor's own construction.

\textbf{Method.} Anchor construction is restricted to sources carrying hard, third-party archival timestamps rather than self-reported metadata. A small human-audited spot-check sample, stratified across year buckets and sources, is then reviewed for signs of misdating, later heavy revision, translation, or ghostwriting, yielding an estimated residual impurity rate. Every downstream false-positive-rate bound is reported net of this impurity estimate, rather than presented as if the anchor's purity were exact.

\subsection{Composition-Matched Resampling}
\textbf{Challenge.} Calendar year and corpus composition --- genre, topic, document length --- co-vary in almost every available continuously-archived corpus, so an apparent false-positive-rate trend over years may reflect composition rather than time.

\textbf{Insight.} The only way to know whether year itself is doing causal work is to build a comparison in which composition cannot vary, and check whether the trend survives when that variation is removed.

\textbf{Method.} A single continuously-archived source is fixed per experiment so that year is the only structural variable, with per-year token budgets and sample sizes matched across buckets. Each year bucket is additionally resampled to equalize its topic-cluster distribution and document-length distribution against a reference bucket, and the drift regression is refit on this composition-matched resample. The observed slope is attributed to genuine temporal change only if it survives --- same sign, comparable magnitude, still significant --- in this matched arm.

\subsection{Reference-Age Propagation Test}
\textbf{Challenge.} A per-document false-positive-rate drift is not automatically informative about the object a downstream user actually consumes --- a corpus-level prevalence number --- because the mixture estimator's mixing coefficient is a nonlinear function of the whole reference distribution, not a simple pass-through of document-level error.

\textbf{Insight.} The only way to measure how much anchor age matters to the reported number is to intervene on the reference directly and observe the corpus-level estimate move, holding everything else about the target corpus fixed.

\textbf{Method.} A target corpus is held fixed --- both a synthetic mixture of known AI-fraction and a real post-2022 corpus of unknown ground truth --- while the human reference bucket fed into the mixture-based prevalence estimator is swapped between a recent bucket and the most temporally distant tested bucket. The resulting shift in the estimated AI-fraction is read off directly as the measured propagation of anchor-age drift into the corpus-level estimand, rather than inferred analytically from the document-level slope.

\subsection{Confidence-Bound Validity Envelope}
\textbf{Challenge.} Per-year sample sizes shrink for older buckets, widening confidence intervals, so drawing a validity boundary at the first year a raw significance test crosses a threshold risks confusing under-power with genuine safety.

\textbf{Insight.} A rate estimate's confidence-interval width is itself information about whether a boundary is certified or merely unmeasured, so building the envelope's edge from the confidence bound rather than the point estimate keeps these two failure modes from being silently conflated.

\textbf{Method.} The validity envelope is defined as the largest contiguous year-range whose upper confidence bound on false-positive rate stays within a pre-declared tolerance of the nominal calibrated rate, at a stated confidence level. This construction is compared directly against a naive envelope drawn at the first year a raw significance test crosses its threshold, to check that the confidence-bound construction is doing genuine work rather than agreeing with the simpler rule by coincidence.

\section{Experiments}

\subsection{Experimental Setup}

\textbf{Datasets.} arXiv abstracts, 2007--2022, drawn from a single continuously-archived source --- the primary pilot source, whose hard timestamps satisfy the purity and leakage constraints directly and whose $\geq$15-year span provides the pre-2018 years needed for the leakage-immune sub-range control. CC-News, a dated news archive --- a scale-up genre with faster topical churn, used to stress-test whether the composition-matching procedure and the drift result generalize beyond formal academic prose. Project Gutenberg, dated by first publication --- extends the anchor's temporal depth well beyond what any recent corpus offers, testing whether the confidence-bound envelope genuinely narrows for very old buckets or merely reflects thinner samples. Reddit/Pushshift comments, post-date stratified --- an informal, high-topic-churn genre used as an additional composition-matching stress test and to diversify the genre coverage the validity envelope must eventually bound. HC3 (Human-ChatGPT Comparison Corpus) --- supplies the AI-side reference needed to build a known-ground-truth mixed target corpus for the reference-age propagation test, where the true AI fraction is known by construction. RAID (Robust AI Detection benchmark) \citep{RAID} --- used exclusively as a pre-registered instrument-validity check: before any threshold is frozen on the calibration split, each instance-level detector is confirmed to reproduce its expected behavior on RAID's known-clean human/AI samples, underwriting the no-threshold-leakage assumption every research question below depends on; it is not itself a claim-bearing dataset for any result.

\textbf{Baselines} (validity-gate and sanity-check comparators, not competitors to beat --- this is a measurement study, and no research question here is framed as ``our procedure beats baseline X''). An anchor-as-asserted control reports current practice directly: assume the nominal calibrated false-positive rate holds uniformly across the whole pre-2022 span; the paper's premise only has stakes if this control is actually violated. A year-shuffle placebo permutes year labels within a source and refits the drift model --- the placebo slope must be null, or the measurement pipeline itself is manufacturing an artifact, a validity gate on the instrument rather than a competitor. The composition-matched resampled arm is the load-bearing rival-explanation control described in the Methodology, not a competitor to beat.

\textbf{Metrics.} Per-bucket false-positive rate and its confidence interval; the drift regression's slope and significance, with and without each disentanglement control applied; the estimated residual anchor-impurity rate; the shift in estimated AI-fraction under reference-year swaps; the boundary year of the validity envelope under both the confidence-bound and naive constructions.

\textbf{Research questions:}
\begin{enumerate}[leftmargin=1.6em]
\item RQ1 --- Does the pre-2022 anchor license a bounded, non-uniform false-positive guarantee at all, as opposed to a flat one or none?
\item RQ2 --- Does the false-positive-rate drift survive restriction to a leakage-immune, pre-LLM sub-range of years?
\item RQ3 --- Does synthetic stylistic perturbation alone reproduce the direction of the observed drift?
\item RQ4 --- How much residual impurity does the certified-human anchor actually carry?
\item RQ5 --- Does the false-positive-rate drift survive composition-matched resampling?
\item RQ6 --- Does anchor age bias the corpus-level prevalence estimate a downstream user actually consumes?
\item RQ7 --- Does a confidence-bound validity envelope draw a meaningfully different boundary than a naive significance-threshold envelope?
\end{enumerate}

\subsubsection{RQ1: Does the pre-2022 anchor license a bounded, non-uniform false-positive guarantee at all, as opposed to a flat one or none?}
\textbf{Setup.} On the primary single-source pilot corpus (arXiv abstracts, $\approx$3,000 documents per year-bucket, thresholds frozen on the 2021 calibration split per Assumption A4), the pooled false-positive rate across all pre-2022 buckets is computed per detector family, together with its 95\% Clopper--Pearson interval, and the per-bucket upper confidence bound on false-positive rate is computed for every year.

\textbf{Expected outcome.} Both hold for at least 2 of the 3 detector families: (i) the pooled anchor false-positive rate is nonzero --- its 95\% Clopper--Pearson lower bound exceeds 0.25\%; and (ii) the bound is not uniform in time --- the maximum per-year upper confidence bound exceeds twice the nominal calibrated rate.

\textbf{Falsification condition.} Refuted if either conjunct fails for a majority of detector families: if the pooled lower bound is indistinguishable from zero, there is no anchor-erosion phenomenon to attribute at all; if the per-year upper confidence bound never exceeds twice the nominal rate, the anchor's guarantee is already uniform across the pre-2022 span and every downstream disentanglement question (RQ2--RQ7) becomes moot rather than merely smaller in magnitude.

\textbf{Hypothesis (plain language).} The pre-2022 anchor's implied error guarantee is real (not simply zero noise) and it is not flat across the years the anchor spans --- establishing that there is a genuine bound-shaped phenomenon for the rest of this study to attribute, propagate, and envelope, rather than presupposing one. Prior fragments make this outcome plausible --- a pooled nonzero pre-LLM false-positive rate for the audited estimator \citep{LLMGeneratedOrHumanWritten} and a simulated year-dependence on the same source \citep{MovingTargetTTA} --- but RQ1 re-establishes it under this design's frozen-threshold, single-source, per-year-bucket protocol rather than importing it from reports whose designs cannot support the attribution questions that follow.

\textbf{Metrics.} Pooled false-positive rate and its Clopper--Pearson interval per detector family; per-bucket upper confidence bound on false-positive rate; the validity envelope $E_k(\epsilon,\delta)$ definition (Problem Definition) evaluated at the pre-registered tolerance.

\textbf{Result.} \ph

\subsubsection{RQ2: Does the false-positive-rate drift survive restriction to a leakage-immune, pre-LLM sub-range of years?}
\textbf{Setup.} Using the primary single-source pilot corpus (arXiv abstracts, $\approx$3,000 documents per year-bucket, with detector thresholds frozen on a 2021 calibration split disjoint from all evaluation years), the drift regression is refit restricted to years at or before 2018 only, for each of three instance-level detector families.

\textbf{Expected outcome.} The slope remains significant and positive --- the same sign as the full-range slope --- for at least 2 of the 3 detector families when restricted to the leakage-immune sub-range.

\textbf{Falsification condition.} Refuted if the slope becomes non-significant or flips sign for a majority of detector families once restricted to years at or before 2018 --- that would mean the full-range slope was being driven by years adjacent to the 2022 cutoff, consistent with undetected leakage rather than a leakage-immune style effect.

\textbf{Hypothesis (plain language).} Whatever drift is observed over the full pre-2022 range is not simply an artifact of years sitting close to the cutoff, where leakage risk is highest.

\textbf{Metrics.} Per-bucket false-positive rate and confidence interval restricted to the sub-range; regression slope and confidence interval on the restricted range.

\textbf{Result.} \ph

\subsubsection{RQ3: Does synthetic stylistic perturbation alone reproduce the direction of the observed drift?}
\textbf{Setup.} A single fixed year's human sample is perturbed at five graded intensities of lexical and syntactic change (synonym substitution rate, sentence-length jitter, punctuation-convention shifts), and each detector's flag rate is measured at each intensity level, using the same fixed sample across all intensities (a paired design).

\textbf{Expected outcome.} Flag rate increases monotonically with perturbation intensity for at least 2 of the 3 detector families, replicating the sign of the slope observed on real year-to-year data, with an absolute flag-rate increase of at least 1 percentage point between the lowest and highest perturbation intensity (the pre-registered H2b effect floor --- significance without magnitude does not pass).

\textbf{Falsification condition.} Refuted if flag rate is flat or non-monotonic across perturbation intensity for a majority of detector families, or if the absolute increase between the lowest and highest intensity falls short of 1 percentage point even where a monotone trend is statistically detectable --- either way, detectors are insensitive (or only negligibly sensitive) to the kind of stylistic change the drift hypothesis requires, undermining the claim that real year-over-year drift reflects a style mechanism at all.

\textbf{Hypothesis (plain language).} Stylistic change of the kind that plausibly occurs over a decade is, by itself, sufficient to move these detectors' flag rate --- establishing that a style mechanism is at least available, independent of whether it is what actually happened historically.

\textbf{Metrics.} Flag rate at each perturbation intensity; trend statistic and significance on flag rate versus intensity.

\textbf{Result.} \ph

\subsubsection{RQ4: How much residual impurity does the certified-human anchor actually carry?}
\textbf{Setup.} A human-audited spot-check sample of at least 2,000 documents, stratified across year buckets and sources, is reviewed for misdating, later heavy revision, translation, or ghostwriting. The sample size is pinned by the pass criterion's own arithmetic, not convenience: under a 95\% Clopper--Pearson bound, an observed impurity of 0.5\% at $n = 1{,}000$ (5 flagged documents) already carries an upper bound of $\approx$1.16\% --- breaching the 1\% cap below even when the expected outcome holds --- whereas at $n = 2{,}000$ the same observed rate bounds at $\approx$0.92\%, so the criterion is decidable at the expected impurity level rather than guaranteed to fail by under-sampling. Per Assumption A1, the audit runs as a validated instrument, not a pass-through metadata read: the rubric is piloted on a $\sim$100-document slice, two independent annotators label the same sample, and inter-rater agreement of $\kappa \geq 0.7$ is required before the impurity estimate is trusted (if $\kappa < 0.7$, the rubric is refined and re-validated before any impurity number is reported).

\textbf{Expected outcome.} The estimated impurity rate stays below 1\%, with its upper confidence bound smaller than the magnitude of any detected drift effect from RQ2/RQ3.

\textbf{Falsification condition.} Refuted if the estimated impurity rate's upper confidence bound exceeds 1\% or is comparable to or larger than the detected drift magnitude --- in that case, any observed false-positive-rate slope cannot be safely attributed to style drift rather than to genuine mislabeled contamination, and every downstream claim in this study would need to be reported net of that uncertainty rather than as a clean drift estimate.

\textbf{Hypothesis (plain language).} The anchor's own authorship certification, while not perfect, is clean enough that its residual uncertainty does not swallow whatever drift effect the study is trying to measure.

\textbf{Metrics.} Annotator-flagged impurity rate and its confidence interval; inter-annotator agreement rate (Cohen's $\kappa$, acceptance gate $\kappa \geq 0.7$).

\textbf{Result.} \ph

\subsubsection{RQ5: Does the false-positive-rate drift survive composition-matched resampling?}
\textbf{Setup.} Each year bucket in the primary single-source corpus is resampled to equalize topic-cluster and document-length distribution against a reference bucket, at matched per-bucket token budgets, and the drift regression is refit on this composition-matched resample.

\textbf{Expected outcome.} The slope survives --- same sign, comparable magnitude, still significant at the pre-registered threshold --- in the composition-matched arm.

\textbf{Falsification condition.} Refuted if the slope disappears or reverses sign once composition is matched --- that would mean the originally observed trend was a genre or topic-composition artifact of corpus assembly, not a temporal style effect, and the stylistic-drift claim specifically would need to be withdrawn even though the anchor's overall invalidity could still stand under a different mechanism.

\textbf{Hypothesis (plain language).} The observed relationship between year and false-positive rate is a property of time itself, not a byproduct of which topics or genres happen to dominate the corpus's older versus newer years.

\textbf{Metrics.} Matched-arm slope and confidence interval; a per-bucket composition-similarity diagnostic relative to the reference bucket.

\textbf{Result.} \ph

\subsubsection{RQ6: Does anchor age bias the corpus-level prevalence estimate a downstream user actually consumes?}
\textbf{Setup.} A target corpus is held fixed --- both a synthetic mixture of known AI-fraction built from anchor human text plus an AI-side reference at a known mixing ratio, and a real post-2022 corpus of unknown ground truth --- while the human reference bucket fed into the mixture-based prevalence estimator is swapped between the most recent tested bucket and the most temporally distant tested bucket.

\textbf{Expected outcome.} The estimated AI-fraction shifts upward by at least 2 percentage points when the aged reference is substituted, holding the target fixed, significant under a paired permutation test.

\textbf{Falsification condition.} Refuted if the estimated AI-fraction is invariant, within its confidence interval, to which reference bucket is used --- meaning the mixture estimator is robust to anchor age, and any per-document drift detected in RQ2/RQ3/RQ5 does not propagate to the corpus-level number end users actually consume.

\textbf{Hypothesis (plain language).} Anchor age does not only move an intermediate detector statistic --- it moves the actual prevalence number a downstream user would report and act on.

\textbf{Metrics.} The shift in estimated AI-fraction between reference-bucket conditions; permutation-test significance; on the known-composition target, the absolute bias against the true mixing ratio.

\textbf{Result.} \ph

\subsubsection{RQ7: Does a confidence-bound validity envelope draw a meaningfully different boundary than a naive significance-threshold envelope?}
\textbf{Setup.} On the same fitted drift curve, the validity envelope is constructed two ways --- a naive construction that draws the boundary at the first year the raw slope's p-value crosses significance, and a confidence-bound construction that draws the boundary at the largest contiguous range whose upper confidence bound on false-positive rate stays within a pre-declared tolerance of the nominal calibrated rate --- and the two boundaries are compared directly.

\textbf{Expected outcome.} The confidence-bound envelope is strictly narrower --- draws the safe boundary at a more recent year --- than the naive envelope, for at least one tested source or genre.

\textbf{Falsification condition.} Refuted if the two constructions draw an identical boundary across all tested sources and genres --- that would mean the confidence-bound construction does no additional work relative to the simpler significance-threshold rule, and the more conservative bookkeeping would be unnecessary complexity rather than a load-bearing contribution.

\textbf{Hypothesis (plain language).} A boundary drawn only from statistical significance is not the same thing as a boundary drawn from genuine certified safety, and the difference between the two is large enough to matter in practice.

\textbf{Metrics.} The boundary year under each construction, per source/genre; the gap, in years, between the two boundaries.

\textbf{Result.} \ph

\vspace{0.5em}
\noindent\textbf{Further evaluation.} With more time, a natural extension is broader language and genre coverage beyond the sources tested here, since the validity envelope this design produces is explicitly source- and genre-scoped rather than a universal claim. A second natural extension, deliberately out of scope for this paper, is proposing a drift-robust prevalence estimator that corrects for the measured bias rather than only reporting it --- that is a repair, not a measurement, and belongs to a different contribution than the one made here.

\section{Related Work}

The closest prior work to this paper is the corpus-level mixture-estimation framework that popularized the pre-2022 anchor \citep{MonitorAIScale}, and it is worth stating the delta plainly before surveying the wider landscape: that framework fits a single mixing coefficient against a fixed human reference distribution and treats it as unproblematic ground truth, but it never tests whether the reference's temporal composition confounds the estimated AI fraction --- which is precisely the instrument this paper audits directly, rather than merely reusing.

\textbf{Corpus-level AI-content prevalence estimation.} The mixture-estimation framework \citep{MonitorAIScale} reframes AI-text detection from instance classification to a population-level measurement problem, fitting a single mixing coefficient between a human-written reference distribution and an AI-generated reference distribution over a stable vocabulary, demonstrated on AI-conference peer reviews. This establishes the corpus-level framing this paper's Results depend on, but treats the human reference corpus as a fixed, unproblematic input rather than an object whose own temporal validity might need auditing --- exactly the gap this paper's reference-age propagation test is built to close. Its own Limitations section explicitly names ``the temporal distribution shift in token frequencies'' as an acknowledged but unmeasured source of error, which this paper turns into the central measured quantity rather than a disclaimed afterthought. A direct follow-on audit \citep{LLMGeneratedOrHumanWritten} applies this estimator to arXiv review and non-review papers and measures its false positives on pre-LLM-era (2020--2022) papers --- 13.1\% and 1.9\% respectively --- then corrects the post-LLM estimates against that pre-LLM baseline. This is the nearest measured evidence for the premise RQ1 tests: the anchor's pooled error is demonstrably nonzero for at least one estimator-corpus pair. But that audit treats the error as a static offset, does not stratify it by year within the pre-LLM span, and states as an explicit limitation that it must assume the false-positive pattern ``remains stable over time'' --- precisely the stability property this paper measures, attributes, and bounds rather than assumes.

\textbf{Temporal concept drift and longitudinal benchmark validation.} A longitudinal malware-detection benchmark spanning more than a decade \citep{LAMDA} holds a consistent feature space fixed across years specifically to measure model degradation over time cleanly, motivated by evidence that naive temporal evaluation inflates performance estimates. This work supplies the year-stratified, temporally-structured validation template this paper adapts, but it targets classification-accuracy decay in a fixed feature space for a different domain (malware), not a calibrated false-positive-rate bound for a text negative-anchor --- its protocol is borrowed, not its findings. A temporal-cutoff contamination diagnostic in financial forecasting \citep{DetectingLookaheadBiasIn} is structurally analogous to this paper's own year-shuffle placebo and pre-LLM sub-range disentanglement --- an interaction test that must vanish in a post-cutoff placebo sample --- confirming the same design pattern generalizes across domains even though its target task (lookahead bias in forecasts) is unrelated to corpus-level AI-content detection.

\textbf{Contamination- and membership-detector assumption fragility.} Work stress-testing benchmark-contamination detectors \citep{BenchContamFragility} shows that the assumption underlying most such detectors --- that seen and unseen samples remain statistically distinguishable --- breaks down under staged post-training, across detector families spanning reference-based, reference-free, perturbation-based, and output-distribution approaches. A reference-free pre-training-data detector \citep{MinKpp} shows a closely related failure mode: membership heuristics degrade toward random as the training and non-training distributions being compared converge. A collection-level contamination detector \citep{ContamICL} further motivates measuring at the corpus level specifically because sample-level signals are confounded by distribution shift. A calibration layer for metric-based detectors \citep{MarkovCalib} shows that detector reliability is brittle under paraphrasing and distribution shift at the mechanism level. All four diagnose that detector reliability collapses as compared distributions converge or shift, but none turns that diagnosis into an explicit false-positive-rate bound as a function of temporal distance for a certified-human anchor --- the temporal gap this paper measures is an unaudited special case of exactly the phenomenon this literature describes only in the abstract.

\textbf{Per-document detection instruments.} A zero-shot detector \citep{DetectGPT} thresholds the local probability-curvature of a passage under a source model, needing no trained classifier or labeled data; a rewrite-based detector \citep{LearnToDistance} instead learns the distance between a passage and its LLM rewrite via normalized log-likelihood rather than a fixed edit or semantic distance; a single-call black-box detector \citep{DnR} shuffles tokens within chunks and uses recovery similarity as a detection proxy, suited to closed-model, low-query deployment. Each of these is validated against contemporaneous human and AI samples, and each is used in this paper's own estimator battery as an instrument --- not a baseline to be beaten, since this is a measurement study, but a component whose false-positive rate on temporally distant human text is exactly what remains uncharacterized in its own original evaluation. A held-out benchmark spanning many models, domains, and adversarial attacks \citep{RAID} is used only to validate that each battery member reproduces its expected behavior before any threshold is frozen, per Assumption A4.

\textbf{Structural exemplars for assumption-auditing measurement studies.} Two recent empirical-study exemplars shape this paper's evidence structure rather than its content: a controlled study interrogating whether larger pretrained language models are actually better calibrated \citep{ACloseLookInto}, and a study measuring whether model-compression techniques degrade fairness across multiple methods and metrics \citep{AComparativeStudyOn}. Both fall on the correct side of the archetype for this paper --- a field-wide operating assumption is turned into a falsifiable, multi-condition measurement rather than argued rhetorically --- and both are imitated here for evidence shape (multi-condition grids, intrinsic-plus-extrinsic metrics, an explicit null/placebo comparison) rather than for content.

\textbf{An alternative, fix-oriented response to temporal drift.} A contemporaneous line of work \citep{MovingTargetTTA} responds to the same underlying observation --- that detector performance degrades under continual distribution shift, including drift in human writing style --- by proposing test-time adaptation: an online, semi-supervised procedure that re-calibrates a detector against unlabeled data as it arrives. Its motivating experiments deserve a plain concession, because they are to date the nearest measured instance of the phenomenon this paper studies: in a simulated deployment calibrated on 2010 arXiv abstracts, a supervised prevalence method estimates 15.1\% AI-generated content in 2020 human abstracts against only 2.3\% in 2010 --- a year-dependent false-positive inflation measured entirely within the pre-ChatGPT span, on the same source this paper's pilot uses. What that measurement is, however, is a motivating observation en route to a repair, not an object of study in its own right: it covers a single detector family on a single source, with no attribution among stylistic change, composition confound, and leakage, no anchor-purity bound, no intervention on the human reference of a mixture-based prevalence estimator, and no validity-envelope construction --- and the repair it motivates is evaluated to a method-superiority standard, not a measurement standard. The two contribution classes remain complementary --- a validated envelope tells a practitioner when reliance on a frozen anchor becomes unsafe; test-time adaptation is one candidate response once it does --- but the honest positioning is that the drift phenomenon has been observed once in passing, while its attribution, propagation, and boundary have not existed before this design.

\section{Conclusion}

This work reframes a widely-used calibration convention --- that pre-2022 human text is a self-certifying, timeless negative anchor for corpus-level AI-content estimators --- as an empirical question with a precise attribution structure, rather than an assumption to take on faith. The central methodological wager is that a false-positive-rate drift is only informative once it survives controls that each rule out a specific rival explanation (composition confound, content leakage, estimator staleness) and its practical stakes are only established once it is shown, not assumed, to move the corpus-level number a downstream estimator actually reports. Whichever outcome the measurement yields, the finding is decision-relevant: a genuine, propagating drift would mean every corpus-level AI-prevalence estimate that reuses this anchor without stating a validity envelope carries an uncharacterized, directional bias that grows as the anchor ages further past 2022; a flat, non-propagating result would instead hand the field its first quantified license for a convention it has so far only assumed. Two limitations bound this claim from the outset and are natural next steps rather than afterthoughts: no scalable oracle can fully certify that a nominally pre-2022 document involved no later revision or mislabeling, so the reported bounds are conditional on the purity spot-check's own coverage, not an absolute guarantee; and the validity envelope this design produces is scoped to the tested sources, genres, and languages, leaving broader-coverage replication --- and any eventual drift-robust estimator repair --- for follow-up work.

\bibliographystyle{plainnat}
\bibliography{refs}

\vspace{1em}
\hrule
\vspace{1em}

\section*{Pre-handoff checklist}
\begin{itemize}[leftmargin=1.4em]
\item \textbf{Organization/navigation} --- narrative flows problem $\to$ method $\to$ experiments $\to$ related work without forward references to undefined terms.
\item \textbf{Terminology discipline} --- ``certified-human anchor,'' ``temporal distance,'' ``validity envelope,'' ``reference-age propagation'' each defined at first use in Problem Definition, used consistently thereafter; no internal stage IDs (\texttt{sp\_*}, \texttt{m\_*}, \texttt{rq\_*}, \texttt{rw\_*}) leaked into prose --- checked directly against the discussion pipeline log.
\item \textbf{Notation with worked example} --- the formal setup section defines $(y^*$, $\Delta(y)$, $A_y$, $\mathrm{FPR}_k(y)$, $\hat{\pi}(T; H))$ and the Problem Definition/Methodology sections consistently reuse this notation; no separate worked numeric example is included since no numbers exist yet (would require fabricating results).
\item \textbf{RQ framing matches outcomes} --- all 7 RQs are graded/quantitative (rates, slopes, boundary years, percentage-point shifts), each answered in its own subsection with an explicit falsifiable expected outcome and a stated falsification condition, not a binary yes/no.
\item \textbf{No typesetting artifacts} --- only explicit placeholder markers in Results, per the drafting convention this is a pre-registered spec, not a report.
\item \textbf{Artifact plan} --- the year-stratified anchor corpus and the estimator-battery/drift-curve code are the natural release artifacts; not yet built, so not claimed as available.
\end{itemize}

\section*{Hand-off}
This is a draft, not a finished paper --- experiments have not run. Natural next steps: prose polish and venue-specific LaTeX formatting once a target venue is chosen, or revising an earlier stage if piloting surfaces a problem (e.g., if the composition-matched arm behaves very differently from the raw arm).

\end{document}
